\documentclass{article}
\usepackage[final]{neurips_2026}

\usepackage[utf8]{inputenc}
\usepackage[T1]{fontenc}
\usepackage{hyperref}
\usepackage{url}
\usepackage{booktabs}
\usepackage{amsfonts}
\usepackage{amsmath}
\usepackage{nicefrac}
\usepackage{microtype}
\usepackage{graphicx}

\title{A scalar forget gate prevents localized storage\\in test-time neural memory}

% Anonymous for double-blind review.
\author{%
  Anonymous Author(s)\\
  Affiliation\\
  \texttt{email}
}

\begin{document}
\maketitle

\begin{abstract}
Test-time neural memory modules are increasingly proposed as the persistent-state
component of long-horizon sequence models, yet little is known about \emph{what}
such a memory stores or \emph{where}. We audit the neural memory of Titans and
report three findings. First, an implementation-fidelity gap: the paper's appendix
specifies a per-dimension decay $\mathrm{diag}(1-\alpha_t)$, but the community
implementation the field builds on emits a single scalar per head, broadcast
identically to every hidden unit. Second, a causal consequence: under the scalar
gate no hidden unit owns any stored association---the largest effect of deleting
any single unit on any tracked pair is $0.13$ cosine---while assigning units
independent decay rates restores localization across $8/8$ seeds, with
single-unit effects $5$--$7\times$ larger. Third, an evaluation trap: the memory's
residual path inflates apparent recall, so that zeroing the entire memory MLP
leaves all but the two freshest associations essentially unchanged. We further
show that whether the memory forgets or accumulates reverses with the training
corpus, and, using a predictability sweep with marginal entropy held fixed,
test whether that reversal tracks temporal structure in the data. Our results
are small-scale but causal, and they bear directly on the reliability of
persistent state in temporal sequence models.
\end{abstract}

\section{Introduction}

Architectures that update a memory \emph{during inference}---test-time training,
neural memory, and state-space hybrids---are being adopted as the persistent-state
mechanism for long-horizon sequence modeling. Titans \citep{behrouz2025titans} is a
prominent instance: a small MLP whose weights are rewritten chunk-by-chunk as the
model reads, governed by a learned forget gate.

The memory update is a discrete-time state-space recursion,
\begin{equation}
M_t = (1-\alpha_t)\,M_{t-1} + S_t,
\label{eq:update}
\end{equation}
where $\alpha_t$ is a decay rate---a \emph{timescale}. If such modules are to serve
as reliable persistent state, we should be able to say what they retain, for how
long, and whether a stored item can be located or corrected. These are the same
questions asked of frozen weights in interpretability research, now posed of a
state that is still moving.

We audit this memory directly, by intervention rather than correlation. Our
contributions are (i) an undocumented divergence between the specification and the
implementation the field actually runs, (ii) a causal demonstration that this
divergence determines whether storage can localize at all, (iii) an evaluation
confound that inflates recall measurements, and (iv) evidence that the memory's
retention regime is set by its training distribution, with a controlled test of
whether temporal predictability is the driver.

\section{Setup}

No official implementation of Titans has been released. We therefore study the
most widely used public reimplementation,\footnote{\texttt{titans-pytorch}, an
unofficial PyTorch implementation. Cited by URL in the camera-ready; omitted here
only where it would identify the authors.} which is the de facto Titans for work
building on the architecture. \textbf{Claims below are properties of that
implementation, not necessarily of the authors' internal code}, which we cannot
inspect.

Unless stated otherwise we use a memory MLP of shape $64\to256\to64$ ($256$ hidden
units) inside a 4-layer Memory-as-Context transformer with $\mathrm{dim}=64$ and a
single memory layer. Two regimes are used: a \emph{toy} regime in which the memory
is trained to recall random key/value vectors, and a \emph{corpus} regime in which
the full model is trained for next-token prediction.

\section{The forget gate is coarser than specified}

Equation~\ref{eq:update} in the main text of \citet{behrouz2025titans} carries no
per-unit index. Appendix~C's fuller derivation writes the decay as
$\mathrm{diag}(1-\alpha_t)$, which would permit each memory dimension its own
retention rate.

In the implementation, the modules producing $\alpha_t$ and the adaptive step size
each emit \textbf{one scalar per head, per chunk}, broadcast identically across all
$256$ hidden units. The finer form in Appendix~C is not available. We found no
statement of this as a deliberate deviation; a cheaper controller and compatibility
with the parallel associative-scan training path are plausible motivations, but we
do not have evidence for either.

This matters because it is a structural constraint, not a tuning choice: with one
shared decay, every unit's retention is coupled, and no unit can specialize in
retaining one item longer than another.

\section{Under the scalar gate, nothing localizes}

\paragraph{Protocol.} We store $N$ tracked key/value pairs in the memory, then
ablate one hidden unit at a time and re-measure the recall of every pair through
the library's own retrieval path, verified to reproduce the unmodified model's
output at $\cos \approx 1.0$.

\paragraph{Result.} No unit's deletion concentrates on any one pair. The largest
effect any single ablated unit had on any single pair's recall was $\mathbf{0.13}$
on a cosine scale, in both trained and untrained memories.

\paragraph{The gate is the cause.} To test whether the shared gate is responsible
rather than some other property of training, we assign each unit its own fixed
decay rate---instantiating Appendix~C's $\mathrm{diag}(1-\alpha_t)$ directly---and
change nothing else. An end-to-end learned per-unit gate did not converge in a
reasonable budget; assigning rates directly isolates the single variable.

\begin{table}[h]
\centering
\caption{Largest single-unit ablation effect on any tracked pair, across
independent seeds. Spread decay exceeds uniform decay's ceiling in every seed.}
\label{tab:perunit}
\begin{tabular}{lcc}
\toprule
Seed & Uniform decay & Per-unit (spread) decay \\
\midrule
0 & 0.019 & 0.157 \\
1 & 0.024 & 0.099 \\
2 & 0.017 & 0.182 \\
3 & 0.015 & 0.158 \\
4--7 & 0.016--0.020 & 0.102--0.207 \\
\midrule
Units cleanly owning one pair & $0/8$ seeds & $4/8$ seeds \\
\bottomrule
\end{tabular}
\end{table}

Localization appears only when units are given independent decay
(Table~\ref{tab:perunit}). We read this as superposition
\citep{elhage2022superposition} in a state written during a single forward pass:
the memory holds more associations than it has independently-decaying slots, so
gradient descent shares units across items. Under the scalar gate the number of
such slots is effectively one.

\paragraph{Editing inherits the entanglement.} Solving for a unit's output row so
that it recalls an arbitrary new target is exact ($\cos = 1.000$ to the new
target); the cost is collateral. Selecting units by ablation-specificity does not
reduce it---a unit roughly $2\times$ more ablation-specific produced \emph{more}
collateral ($0.540$ vs.\ $0.473$), because ablation-specificity measures how small
a unit's existing contribution to other pairs is, while editing collateral depends
on how strongly it activates for their keys.

\section{An evaluation trap: the residual floor}

The default memory is $\mathrm{ResidualNorm}(\mathrm{MemoryMLP}(x))$, i.e.\
$\mathrm{LN}(\mathrm{MLP}(q))\cdot\gamma + q$. Zeroing the \emph{entire} MLP---
removing all learned memory content---left most pairs' recall nearly unchanged;
only the two freshest showed a real contribution ($0.96\to0.59$ and
$0.97\to0.57$). Most of what a naive read scores as ``recall'' is the residual
path returning the query, not stored content.

Any evaluation of this architecture that does not zero the memory as a control
will overstate retention, and the overstatement is largest for exactly the old
items long-horizon claims depend on.

\section{Retention regime is set by the training distribution}

In the toy regime the memory forgets sharply: the first write decays to a few
percent of peak within six chunks (end/early write-norm ratio $0.05$--$0.06$). On
natural text the same measurement \emph{reverses}: across five independent runs the
ratio is $1.42$--$2.21$, i.e.\ the memory accumulates, with the effect strongest in
the better-trained runs.

We tested two mechanisms and ruled both out. Hooking the decay module directly and
comparing its output on real text against random bytes through identical weights
showed no stable gap ($\le 0.033$, and in the direction opposite the hypothesis).
Measuring $\cos(\text{write}_t, \text{write}_{t-1})$ showed no stable difference
between real and random input. What remains is that the regime is fixed by what the
memory was \emph{trained} on, rather than decided live from what it reads.

\paragraph{A controlled test.} To test that directly we sweep temporal
predictability while holding the marginal distribution fixed. Each corpus---AR(1)
processes with increasing $\phi$, a periodic signal, a real forecasting benchmark,
and text---is mapped to $256$ equal-frequency bins, so all corpora have identical
marginal entropy ($5.545$ nats) and differ only in temporal dependence
(lag-1 mutual information $0.002$ to $1.71$ nats). The same architecture is trained
on each, and the write-norm ratio is measured on held-out data.
\textbf{[Results pending; see Table~\ref{tab:temporal}.]}

\begin{table}[h]
\centering
\caption{Retention regime versus temporal predictability, marginal entropy matched.}
\label{tab:temporal}
\begin{tabular}{lccc}
\toprule
Corpus & Lag-1 MI (nats) & Val.\ loss (nats) & Write-norm ratio \\
\midrule
AR(1), $\phi=0.00$ & 0.002 & --- & --- \\
AR(1), $\phi=0.50$ & 0.138 & --- & --- \\
Periodic & 0.686 & --- & --- \\
AR(1), $\phi=0.90$ & 0.778 & --- & --- \\
AR(1), $\phi=0.99$ & 1.706 & --- & --- \\
Real benchmark & --- & --- & --- \\
Text & --- & --- & --- \\
\bottomrule
\end{tabular}
\end{table}

\section{Related work}

\citet{titansrevisited2025} reimplement Titans and report that chunking limits
performance and that memory updates alone do not deliver meaningful test-time
learning. That work evaluates behavior at the benchmark level; ours opens the
memory and asks where content sits, so the two are complementary.

\citet{young2026writesae} learns rank-1 dictionary atoms over the writes of Gated
DeltaNet, Mamba-2, and RWKV-7, and reports both localized and superposed atoms with
little collateral from erasure. We study a different architecture with a different
method---direct causal ablation of raw units, with no learned decomposition---and
find substantial collateral. We take this as consistent rather than contradictory:
dictionary learning is designed to recover a sparse, reusable basis, which is
precisely the entanglement that raw units exhibit. Applying such a decomposition to
Titans before editing is a natural next step.

\citet{yildirim2026superposition} finds, using sparse autoencoders on a forecasting
transformer, that superposition is \emph{not} necessary for competitive time-series
forecasting. Our setting differs in that the state under study is written at test
time, and we find entanglement to be forced by the gate rather than optional.

\section{Limitations}

All results are small-scale: a $256$-unit memory, short documents, and short
training runs. We study one architecture and one public implementation, and cannot
verify whether the authors' unreleased code shares the scalar-gate restriction. The
per-unit decay result uses directly assigned rates rather than a learned gate, so
it establishes that per-unit decay \emph{permits} localization, not that a trained
model would discover it. The cause of the retention reversal remains open; the
sweep in Section~6 narrows it but does not by itself identify a mechanism.

\begin{thebibliography}{9}
\small
\bibitem[Behrouz et al.(2025)]{behrouz2025titans}
A.~Behrouz, P.~Zhong, and V.~Mirrokni.
\newblock Titans: Learning to memorize at test time.
\newblock \emph{NeurIPS}, 2025. arXiv:2501.00663.

\bibitem[Elhage et al.(2022)]{elhage2022superposition}
N.~Elhage et al.
\newblock Toy models of superposition.
\newblock \emph{Transformer Circuits Thread}, 2022. arXiv:2209.10652.

\bibitem[Titans Revisited(2025)]{titansrevisited2025}
Titans revisited: A lightweight reimplementation and critical analysis of a
test-time memory model.
\newblock 2025. arXiv:2510.09551.

\bibitem[Young(2026)]{young2026writesae}
J.~Young.
\newblock WriteSAE: Sparse autoencoders for recurrent state.
\newblock 2026. arXiv:2605.12770.

\bibitem[Y{\i}ld{\i}r{\i}m(2026)]{yildirim2026superposition}
A.~Y{\i}ld{\i}r{\i}m.
\newblock Superposition is not necessary: A mechanistic interpretability analysis
of transformer representations for time series forecasting.
\newblock 2026. arXiv:2605.05151.
\end{thebibliography}

\end{document}
